\section{Brainstorming（头脑风暴）}

\subsection{目标与原则}

Brainstorming 解决的是 \textbf{What \& Why}——"我们要建什么、为什么这样建"。

\begin{itemize}
    \item 一次只问一个确认性问题。
    \item 理解用户的真实意图（Purpose）、约束和成功准则。
    \item 先探索上下文，提出 \textbf{2--3 个方案}，分析各自的 Trade-off。
    \item 用户确认后，形成并保存 \textbf{Design Doc}（\texttt{docs/plans/<date>-<topic>-design.md}）。
\end{itemize}

\begin{warningbox}{头脑风暴阶段的禁令}
在 Brainstorming 阶段，\textbf{禁止调用任何实现工具、禁止写任何代码}。此阶段的产出物纯粹是概念与决策文档，侧重宏观架构、数据流向、组件交互、异常处理。
\end{warningbox}

\subsection{流程步骤}

\begin{enumerate}
    \item 探索项目上下文。
    \item 向用户提出确认性问题，明确问题边界。
    \item 提出 2--3 个解决方案，分析 Trade-off。
    \item 呈现 Design Section，等待用户批准。
    \item 用户批准后，自动调用 \texttt{Writing Plans} Skill。
\end{enumerate}

\subsection{本章小结}

头脑风暴是需求澄清的必经阶段，产出 Design Doc，为后续实施计划提供架构基础。其核心纪律是"只设计不编码"。


